%%%%%%%%%%%%%%%%%%%%%%%%%%%%%%%%%%%%%%%%%%%%%%%%%%%%%%%%%%%%%%
%% Chapter 1: Introduction
%%%%%%%%%%%%%%%%%%%%%%%%%%%%%%%%%%%%%%%%%%%%%%%%%%%%%%%%%%%%%%

\section{Introduction}
\label{chap:introduction}

\subsection{Motivation and context}
\label{sec:motivation-context}

Sign languages are fully-fledged natural languages, possessing their own lexicons,
grammars, and morphosyntactic structures that are independent of the spoken
languages with which they coexist. Contrary to common assumption, they are neither
universal nor manual transliterations of speech: American Sign Language (ASL),
British Sign Language (BSL), German Sign Language (DGS), and Polish Sign Language
(PJM) are mutually unintelligible, and the grammar of a sign language frequently
bears no resemblance to that of the surrounding spoken language. ASL, for instance,
does not follow English word order, and conveys grammatical information --
negation, interrogation, conditionality -- through non-manual markers such as
facial expression, head position, and eye gaze, in parallel with the manual
articulation of signs. For the Deaf and hard-of-hearing communities for whom a
sign language is the primary language, this independence creates a genuine
communication barrier with the hearing majority, and motivates the development of
automatic sign language translation systems.

Automatic sign language translation is not a single task but a family of
interrelated sub-problems. Translation between a spoken language and a sign
language may proceed in either direction, and is conventionally decomposed around
an intermediate symbolic representation known as a \emph{gloss}. A gloss is a
written transcription of a sign sequence, typically rendered as capitalised words
of a spoken language, that records which signs are produced and in what order
while abstracting away their visual realisation. For example, the DGS utterance
glossed as ICH OSTERN WETTER ZUFRIEDEN corresponds to the fluent German sentence
``Da k\"onnen wir mit unserem Osterwetter eigentlich ganz zufrieden sein'' --
demonstrating that the gloss is a compact, lossy intermediate rather than a
complete translation. This decomposition yields four directed sub-tasks:
video-to-gloss and gloss-to-text on the recognition side, and text-to-gloss and
avatar-based gloss-to-video on the production side. These sub-tasks, and how the
present project positions itself among them, are discussed in more detail in
Chapter~\ref{chap:related-work}.

In view of the limited annotated data resources available for sign language
research, and the practical need for tooling to collect and label such data, this
project combines a research track and a development track. The research track
targets the video-to-gloss recognition problem and its extension towards
continuous, sentence-level translation. The development track builds the
supporting software infrastructure -- a data annotation platform and an
educational demonstration platform -- required to collect training data and to
present the resulting models to end users.

\subsection{Goals and scope}
\label{sec:goals-scope}

The main objective of this project is to research, develop, and implement methods
for translating sign language into written or spoken language, and to build the
surrounding software systems that make this research reproducible and usable. In
view of the limited data resources available on the market, the project
concentrates on both the research problem itself and on producing a comprehensive
system with tangible practical and scientific value. Concretely, the project
pursues the following goals:

\begin{itemize}
  \item design and implement a decentralized data annotation platform for
        collecting and labelling sign language video datasets;
  \item formulate and evaluate a video-to-gloss recognition model for isolated
        sign gestures, benchmarking multiple temporal encoder architectures under
        identical experimental conditions;
  \item investigate translation between gloss sequences and natural-language
        text in both directions (gloss-to-text and text-to-gloss);
  \item integrate the isolated recognition and translation components into an
        end-to-end, sentence-level video-to-text pipeline;
  \item deliver an educational platform that demonstrates the trained models to
        end users through an accessible interface.
\end{itemize}

\todo{PLACEHOLDER: State the precise scope boundaries of the thesis -- which sign
language(s) and dataset(s) are targeted, what is explicitly out of scope (e.g.
continuous/real-time translation, multi-signer deployment, avatar-based
production), and any constraints inherited from the source data (e.g. isolated,
pre-segmented gestures only).}

\subsection{Thesis structure}
\label{sec:thesis-structure}

The remainder of this thesis is organised as follows.
Chapter~\ref{chap:related-work} surveys prior work in sign language recognition
and translation, covering video-to-gloss methods, NLP models applied to sign
language, and existing systems and benchmarks.
Chapter~\ref{chap:annotation-platform} describes the annotation platform built to
collect and label the video datasets used throughout this work.
Chapter~\ref{chap:isolated-gloss-recognition} presents the isolated gloss
recognition (video-to-gloss) model, including the dataset, preprocessing
pipeline, model architecture, and experimental results.
Chapter~\ref{chap:text-translation} addresses translation between gloss
sequences and natural-language text.
Chapter~\ref{chap:sentence-level-translation} integrates the components from
Chapters~\ref{chap:isolated-gloss-recognition} and~\ref{chap:text-translation}
into an end-to-end, sentence-level video-to-text pipeline.
Chapter~\ref{chap:educational-platform} presents the educational platform that
demonstrates the resulting models to end users.
Finally, Chapter~\ref{chap:conclusions} answers the goals set out above,
discusses the limitations of the present system, and outlines directions for
future work.
